% wave13_a4_weyl_kac_product_polyakov.tex
% Polyakov-voice adversarial assault on the Weyl-Kac-Borcherds denominator
% identity for g_{Delta_5}, and on the infinite product expansion of
% Delta_5 as the Borcherds singular theta lift of phi_{0,1}.
%
% Channeling: Polyakov bootstrap rigour + Russian math physics elite
% school (Gelfand-Manin-Drinfeld-Arnold-Beilinson + Nekrasov-Witten-
% Costello-Gaiotto-Moore). Costello-Francis-Gwilliam chain-level
% elementarity in each Heal.
%
% Target primary sources (paper transcript, Lorgat 2020):
%   Section 2  -- Delta_5 preamble, f(n,l,m), triple-product seed
%                 identity 1 + (1/64) sum f(1+2t,1,1) q^t = prod(1-q^k)^9
%   Section 6  -- phi_{12,1} / delta_{12} = phi_{0,1}, weak Jacobi form
%                 weight 0 index 1; Theorem 4 infinite product identity
%                 for Delta_5 with exponents f(nm, l)
%   plus the Borcherds 1998 weight theorem 13.3:
%                 wt(Phi) = c_0(0)/2
%                 applied with c_1(0) = f_{phi_{0,1}}(0,0) = 10.
%
% Mission: >=5 ATTACK <-> HEAL cycles on
%   (i)   phi_{0,1} q-expansion + theta decomposition + Fourier
%         coefficients c(n,r) with 4n - r^2 >= -1;
%   (ii)  Borcherds singular theta lift
%           sum c(n,r) -> prod (1 - p^a q^b zeta^c)^{c(ab,c)}
%         over (a,b,c) > 0 in the positive Weyl chamber of Lambda^{2,1}_{II};
%   (iii) convergence to Delta_5 as weight-5 form on Sp_4(Z);
%   (iv)  Weyl-Kac-Borcherds denominator formula for g_{Delta_5}
%           e^rho prod_{alpha>0} (1 - e^{-alpha})^{mult(alpha)};
%   (v)   super-dim of root spaces matches phi_{0,1} Fourier data via
%         Gritsenko-Nikulin duality.
%
% Authorship: Raeez Lorgat. Scope: notes/ (not manuscript).

\documentclass[11pt]{article}
\usepackage[margin=1in]{geometry}
\usepackage{amsmath,amssymb,amsthm,mathtools}
\usepackage{enumitem}
\usepackage{hyperref}

\newtheorem{theorem}{Theorem}
\newtheorem{proposition}[theorem]{Proposition}
\newtheorem{lemma}[theorem]{Lemma}
\newtheorem{corollary}[theorem]{Corollary}
\theoremstyle{definition}
\newtheorem{definition}[theorem]{Definition}
\newtheorem{construction}[theorem]{Construction}
\theoremstyle{remark}
\newtheorem{attack}{Attack}
\newtheorem{heal}{Heal}
\newtheorem{observation}[theorem]{Observation}
\newtheorem{warning}[theorem]{Warning}
\newtheorem{verification}[theorem]{Verification}
\newtheorem{ghost}[theorem]{Ghost theorem}

\newcommand{\CC}{\mathbb{C}}
\newcommand{\ZZ}{\mathbb{Z}}
\newcommand{\QQ}{\mathbb{Q}}
\newcommand{\RR}{\mathbb{R}}
\newcommand{\NN}{\mathbb{N}}
\newcommand{\HH}{\mathbb{H}}
\newcommand{\Lat}{\Lambda}
\newcommand{\fg}{\mathfrak{g}}
\newcommand{\fn}{\mathfrak{n}}
\newcommand{\fh}{\mathfrak{h}}
\newcommand{\Sp}{\mathrm{Sp}}
\newcommand{\OO}{\mathrm{O}}
\newcommand{\SO}{\mathrm{SO}}
\newcommand{\GL}{\mathrm{GL}}
\newcommand{\SL}{\mathrm{SL}}
\newcommand{\Wg}{\mathsf{W}}
\newcommand{\sgn}{\operatorname{sgn}}
\newcommand{\mult}{\operatorname{mult}}
\newcommand{\re}{\mathrm{re}}
\newcommand{\im}{\mathrm{im}}
\newcommand{\ch}{\mathrm{ch}}
\newcommand{\Borch}{\mathrm{Borch}}
\newcommand{\Lift}{\mathrm{Lift}}
\newcommand{\Delta}{\varDelta}  % use \varDelta for Delta symbol; discard Greek conflict
\renewcommand{\Delta}{\mbox{\ensuremath{\Delta}}}

\title{Weyl--Kac product for $\fg_{\Delta_5}$ and the Borcherds lift of
$\phi_{0,1}$:\\ Polyakov bootstrap attack--heal}
\author{Raeez Lorgat}
\date{2026-04-22}

\begin{document}
\maketitle

\begin{abstract}
Target: the pair of identities (Thm.~3, Thm.~4 of Lorgat~2020):
\[
  \tfrac{1}{64}\Delta_5(2Z)
   \;=\; \sum_{w\in \Wg^{(2)}(\Lat^{2,1}_{II})}\!\!\det(w)
         \Bigl(e^{-\pi i(w\rho,z)}
          - \!\!\sum_{a\in\Lat^{2,1}_{II}\cap\RR_{>0}\mathcal P_{II}}\!\!
            m(a)\,e^{-\pi i(w(\rho+a),z)}\Bigr),
\]
\[
  \tfrac{1}{64}\Delta_5(Z)
   \;=\; e^{\pi i(z_1+z_2+z_3)}
         \prod_{(n,\ell,m)>0}(1 - e^{2\pi i(n z_1+\ell z_2+m z_3)})^{f(nm,\ell)},
\]
the second being the infinite-product expansion witnessing $\Delta_5$ as
a Borcherds singular theta lift of $\phi_{0,1}$, the first the
Weyl--Kac--Borcherds denominator formula for the generalized
Borcherds--Kac--Moody Lie superalgebra $\fg_{\Delta_5}$ associated to
the Gram matrix $((\delta_i,\delta_j))_{i,j=1,2,3}$ with off-diagonal
$-2$ and diagonal $2$. The two are the same formal identity read in
different lanes: Cartan--Eilenberg side (Weyl--Kac) and
Selberg--Borcherds theta-integral side.

Five attack$\leftrightarrow$heal cycles probe each load-bearing step:
(1)~the Eichler--Zagier theta decomposition of $\phi_{0,1}$ and
the Fourier data $c(n,r)$; (2)~the Borcherds singular theta lift and
the positive-Weyl-chamber exponent bookkeeping; (3)~weight
$=c_1(0)/2=5$ via Borcherds 1998 Thm.~13.3; (4)~Weyl--Kac--Borcherds
denominator formula from the GBKM generators and Serre relations;
(5)~Gritsenko--Nikulin duality
$\mult(\alpha_{(a,b,c)})=f(ac,b)=c_{\phi_{0,1}}(ac,b)$ relating root
multiplicities to Jacobi Fourier coefficients. Each cycle is a
Polyakov-bootstrap probe: a conformal-invariance-style constraint plus
a worked counter-example hammer. Claim status honest throughout.
\end{abstract}

\tableofcontents

%====================================================================
\section{Load-bearing primary data}
\label{sec:data}

Recall the three lattices and the Weyl vector from Lorgat~2020
\S3--\S4.

\begin{itemize}[leftmargin=1.5em]
\item $\Lat^{3,2}\simeq\Lat^{1,1}\oplus\Lat^{1,1}\oplus[2]$: signature
$(3,2)$, hyperbolic-$\oplus$-hyperbolic-$\oplus$-rank-$1$ with
Gram $(2)$.
\item $\Lat^{2,1}=\Lat^{1,1}\oplus[2]\simeq\ZZ f_2\oplus\ZZ f_3\oplus
\ZZ f_{-2}\subset\Lat^{3,2}$ (primitive hyperbolic sublattice, signature
$(2,1)$).
\item $\Lat^{2,1}_{II}=\{m f_2+\ell f_3+n f_{-2}\in\Lat^{2,1}\mid
m\equiv n\equiv 0\bmod 2\}$, the type-II sublattice of index $6$ in
$\Lat^{2,1}$, with fundamental polyhedron $\mathcal P_{II}$ cut by
$\{\delta_1=2f_2-f_3,\ \delta_2=2f_{-2}-f_3,\ \delta_3=f_3\}$. Gram:
\begin{equation}
\label{eq:gram-delta-ij}
  (\delta_i,\delta_j)_{i,j=1,2,3}\;=\;
  \begin{pmatrix} 2 & -2 & -2\\ -2 & 2 & -2\\ -2 & -2 & 2 \end{pmatrix}.
\end{equation}
\item Weyl vector $\rho=\tfrac12\delta_1+\tfrac12\delta_2+\tfrac12\delta_3
=f_2-\tfrac12 f_3+f_{-2}$; satisfies $(\rho,\delta_i)=-1=-(\delta_i,
\delta_i)/2$ for $i=1,2,3$.
\item $\mathcal P_{II}\subset\mathcal C(\Lat^{2,1})_+/\RR_{>0}$ is
finite-volume in the hyperbolic plane; three vertices at infinity
$\{2f_2,\ 2f_{-2},\ 2f_2-2f_3+2f_{-2}\}$ sit on the null cone
$\{a\in\Lat^{2,1}_{II}\mid(a,a)=0\}\cap\RR_{>0}\mathcal P_{II}$, each
orbit-related under $\mathrm{Aut}(\mathcal P_{II})\simeq S_3$.
\item $\phi_{0,1}\in J^{\mathrm{weak}}_{0,1}(\SL_2(\ZZ))$: weight $0$,
index $1$, weak Jacobi form defined by $\phi_{0,1}=\phi_{12,1}/\delta_{12}$
with $\delta_{12}=\eta^{24}$ and $\phi_{12,1}$ the first Fourier--Jacobi
coefficient of $\Delta_{12}\in M_{12}(\Sp_4(\ZZ))$. Fourier expansion
(Eichler--Zagier 1985 Thm.~2.2, Lorgat~2020 \S6):
\begin{equation}
\label{eq:phi01-expansion}
  \phi_{0,1}(\tau,z)\;=\;\sum_{n\geq 0,\ell\in\ZZ}f(n,\ell)\,q^n r^\ell,
  \qquad q=e^{2\pi i\tau},\ r=e^{2\pi i z},
\end{equation}
with $f(n,\ell)$ depending only on the discriminant $4n-\ell^2$, polar
range $4n-\ell^2\geq -1$, and values
\[
  f(0,0)=10,\quad f(0,\pm 1)=1,\quad f(0,\ell)=0\text{ for }|\ell|\geq 2,
\]
\[
  f(1,0)=-64,\quad f(1,\pm 1)=108,\quad f(1,\pm 2)=-64,\quad f(1,\pm 3)=10,\ldots
\]
Equivalently the $q^0$ Laurent polynomial is
$\phi_{0,1}|_{q=0}=r^{-1}+10+r$.
\item The Igusa form $\Delta_5\in M_5(\Sp_4(\ZZ),\nu_{\Delta_5})$ of
weight $5$ with order-$2$ multiplier $\nu_{\Delta_5}$ (Maass~1964);
Fourier expansion in $Z=\begin{psmallmatrix}z_0&z_1\\z_2&z_3
\end{psmallmatrix}$ has $f(1,1,1)=64$, $64\mid f(n,\ell,m)$.
\end{itemize}

\begin{observation}[The $f$ of \S2 and the $f$ of \S6 are
\emph{the same} Fourier coefficients]
\label{obs:f-f-identification}
Lorgat~2020 writes $f(n,\ell,m)$ for the three-variable Fourier
coefficients of $\Delta_5(Z)$ (\S2 eq.~following Maass, p.~3 line 174),
and $f(n,\ell)$ for the two-variable Fourier coefficients of
$\phi_{0,1}(\tau,z)$ (\S6 p.~9 line 717). The notation collides; what
Theorem~4 says is the \emph{non-trivial} identification
\[
  f(n,\ell,m)\big|_{\text{via Borcherds lift}}
  \;=\; f(nm,\ell)\big|_{\text{Jacobi form}},
\]
i.e.\ the Fourier coefficient of $\Delta_5$ at multi-index $(n,\ell,m)$
equals the Fourier coefficient of $\phi_{0,1}$ at discriminant
$(nm,\ell)$. This is the content of the infinite product. The
$q^0$-slice identification $r^{-1}+10+r\rightsquigarrow
c_1(0)=10=\sum\chi(K3)^{\mathrm{ell,sym}}_{\mathrm{massless}}$ feeds
the Borcherds weight theorem.
\end{observation}

%====================================================================
\section{Cycle 1: Theta decomposition of $\phi_{0,1}$ and the
$c(n,r)$ data}
\label{sec:cycle1}

\begin{attack}[Which decomposition, which index, what convention?]
\label{atk:cycle1}
Lorgat~2020 quotes Theorem~4
\[
  \tfrac{1}{64}\Delta_5(Z)
  \;=\; e^{\pi i(z_1+z_2+z_3)}\prod_{(n,\ell,m)>0}
        (1-e^{2\pi i(n z_1+\ell z_2+m z_3)})^{f(nm,\ell)}
\]
and labels $f(nm,\ell)$ the Fourier coefficient of $\phi_{0,1}$ at
discriminant $(nm,\ell)$. Four objections before the lift is accepted.

\begin{enumerate}[label=\textbf{O.\arabic*},leftmargin=1.5em]
\item The Borcherds singular theta lift takes as input a
\emph{vector-valued} modular form of weight $(2-\mathrm{rk}\Lat)/2=1/2$
on $\mathrm{Mp}_2(\ZZ)$ with values in the Weil representation of
$\Lat^{3,2}$ (Borcherds 1998 \S3--\S4). The starting $\phi_{0,1}$ has
weight $0$ and index $1$, not $1/2$ with vector values. A conversion
from weak Jacobi form to Weil-representation-valued modular form must
intervene, and is never stated in Lorgat~2020.
\item The exponent $f(nm,\ell)$ is indexed by $(nm,\ell)$, not by
$(n,\ell,m)$ directly. Why should the lift only care about the product
$nm$? This is a non-trivial Eichler--Zagier theta-decomposition
phenomenon and warrants a first-principles derivation, not a citation.
\item The positivity condition ``$(n,\ell,m)>0$'' is defined in
Remark~1 (Lorgat 2020 p.~9 line 735) as ``$n\geq 0,m\geq 0,\ell$
arbitrary if $n>0$ or $m>0$, and $\ell<0$ if $n=m=0$''. This is the
positive Weyl chamber of $\Lat^{2,1}_{II}$ with respect to a chosen
Weyl vector. But the Weyl vector $\rho=f_2-\tfrac12f_3+f_{-2}$ uses
\emph{half-integer} coefficients, while the sum is indexed by
integer $(n,\ell,m)$. How does this half-integer Weyl vector reconcile
with integer lattice points?
\item Convergence of the product is asserted via
$f(n,\ell)=O(\exp(\sqrt{4n-\ell^2}))$, an asymptotic statement. But
\emph{convergence} requires a more delicate bound in a neighbourhood
of the zero-dimensional cusp (the standard Hecke-type estimate is not
enough); the argument passes to ``methodology of Kac [11]'' without
displaying the actual convergence inequality.
\end{enumerate}
\end{attack}

\begin{ghost}[Eichler--Zagier theta decomposition as Weil-rep conversion]
\label{ghost:cycle1}
Let $\phi\in J^{\mathrm{weak}}_{k,m}(\SL_2(\ZZ))$ be a weak Jacobi form
of integer weight $k$ and integer index $m$. The Eichler--Zagier
theta decomposition (Eichler--Zagier 1985 Thm.~5.1; cf.\
Bruinier~2002 Chap.~1.2) asserts
\begin{equation}
\label{eq:theta-decomp}
  \phi(\tau,z)\;=\;\sum_{\mu\in\ZZ/2m\ZZ}h_\mu(\tau)\,
  \vartheta_{m,\mu}(\tau,z),
\end{equation}
where $\vartheta_{m,\mu}(\tau,z)=\sum_{\nu\in\ZZ+\mu/2m}q^{m\nu^2}r^{2m\nu}$
is the Jacobi theta, and the vector $(h_\mu)$ transforms as a
\emph{vector-valued} modular form of weight $k-1/2$ on the metaplectic
cover $\mathrm{Mp}_2(\ZZ)$ under the Weil representation
$\rho_m$ associated to the rank-$1$ positive-definite lattice $[2m]$.

For $\phi=\phi_{0,1}$: $k=0$, $m=1$, hence $(h_\mu)$ is weight $-1/2$,
Weil representation of $[2]$. The Fourier expansion
$h_\mu(\tau)=\sum_{D}C(D)q^{D/4m}$ (where $D=4mn-\ell^2$ with
$\ell\equiv\mu\bmod 2m$) packages the $f(n,\ell)=C_\mu(4n-\ell^2)$
into \emph{discriminant-indexed} Fourier coefficients. This explains
Objection O.1: the conversion is the theta decomposition; $(h_\mu)$ is
the Weil-representation-valued input.

For $m=1$: two components $\mu=0,1$ with $h_\mu(\tau)\in M^!_{-1/2}
(\rho_1)$. The principal part (polar discriminants $D\leq 0$) fits in
the set $D\in\{-1,0\}$:
\[
  C_1(-1)=1\quad(\text{from }f(0,\pm 1)=1),\qquad
  C_0(0)=10\quad(\text{from }f(0,0)=10).
\]
\emph{This is the input to Borcherds's singular theta lift.} The entry
$C_0(0)=10$ is the $c_1(0)$ that enters the weight formula
$\mathrm{wt}(\Phi)=c_0(0)/2=5$.

Objection O.2 resolution: the exponent $f(nm,\ell)$ depends on
$(nm,\ell)$ alone because after theta decomposition, the
Borcherds--Gritsenko identity maps \emph{each pair} $(4nm-\ell^2,
\ell\bmod 2)$ to a Fourier coefficient $C_\mu(D)$ at $D=4nm-\ell^2$,
$\mu=\ell\bmod 2$. The two variables $(n,m)$ combine only via the
discriminant $D=4nm-\ell^2$, which is the symplectic-invariant
combination under the swap $n\leftrightarrow m$
(the involution $Z\leftrightarrow Z^t$ inside $\Sp_4(\ZZ)$).

Objection O.3 resolution: the Weyl vector $\rho=\tfrac12\delta_1+
\tfrac12\delta_2+\tfrac12\delta_3=f_2-\tfrac12f_3+f_{-2}$ sits in
$(\Lat^{2,1}_{II})^*=\tfrac12\Lat^{2,1}$, not in $\Lat^{2,1}_{II}$
itself. The exponent of $e^{\pi i(\rho,z)}$ in the product side of
Thm.~4 is $\tfrac12(\rho+\rho')$ for the \emph{rescaled} Weyl vector
$\rho'=f_2+f_3+f_{-2}$ matching the $(n,\ell,m)$-indexing
(cf.\ $e^{\pi i(z_1+z_2+z_3)}$ in Thm.~4: the coefficients are the
integer triple $(1,1,1)$, which is $2\rho$ under the pairing with the
basis $\{f_{-2},f_3,f_2\}$). Integer-vs-half-integer is a change of
basis, not a contradiction.

Objection O.4 resolution: the asymptotic $f(n,\ell)=O(\exp(\pi
\sqrt{4n-\ell^2}))$ is the Hardy--Ramanujan--Rademacher estimate
for Jacobi forms of weight $k\leq 0$ (Dabholkar--Murthy--Zagier 2012
Lemma~8.1; Bringmann--Folsom--Ono--Rolen 2017 Thm.~7.5). Convergence
of the product on neighbourhoods of the zero-dimensional cusp then
follows from Borcherds 1998 Thm.~10.1 (\emph{not} 13.3): the standard
cusp-neighbourhood convergence is the content of Borcherds's Theorem~7.3
plus the asymptotic. The exponential weight
$\exp(-2\pi(n y_1+\ell y_2+m y_3))$ with $y_i=\mathrm{Im}(z_i)$ combined
with the polynomial growth of $f(nm,\ell)$ in $\sqrt{4nm-\ell^2}$
gives absolute convergence on
$\{Z\mid y_1 y_3-y_2^2\geq C,\ y_1,y_3\geq C'\}$ for any $C,C'>0$.
\end{ghost}

\begin{heal}[First-principles restatement of the theta decomposition]
\label{hl:cycle1}
Inscribe \eqref{eq:theta-decomp} with specific values for $\phi_{0,1}$
and identify the Weil-representation input explicitly.

\textbf{Step~1 (Theta components).} Write
\[
  \phi_{0,1}(\tau,z)\;=\;h_0(\tau)\vartheta_{1,0}(\tau,z)
  +h_1(\tau)\vartheta_{1,1}(\tau,z),
\]
with
\[
  \vartheta_{1,0}(\tau,z)=\sum_{\nu\in\ZZ}q^{\nu^2}r^{2\nu},\qquad
  \vartheta_{1,1}(\tau,z)=\sum_{\nu\in\ZZ+1/2}q^{\nu^2}r^{2\nu}.
\]
Reading off the $r^\ell$ coefficients of $\phi_{0,1}$ at each $q^n$:
\[
  h_0(\tau)\;=\;\sum_{\substack{D\geq -0\\D\equiv 0\bmod 4}}c_0(D)q^{D/4},
  \qquad
  h_1(\tau)\;=\;\sum_{\substack{D\geq -1\\D\equiv -1\bmod 4}}c_1(D)q^{D/4},
\]
with principal parts
\[
  c_0(0)=10,\qquad c_1(-1)=1,
\]
and higher Fourier coefficients $c_\mu(D)=f(n,\ell)$ for $D=4n-\ell^2$,
$\ell\equiv\mu\bmod 2$.

\textbf{Step~2 (Weil-representation reading).} The pair $(h_0,h_1)$
transforms as a vector-valued modular form of weight $-1/2$ on
$\mathrm{Mp}_2(\ZZ)$ under the Weil representation $\rho$ of the
rank-$1$ positive-definite lattice $[2]\simeq\ZZ\cdot\delta_3$
(Eichler--Zagier 1985 Thm.~5.1; Skoruppa 1990 Lemma~2.2):
\[
  (h_0,h_1)\bigl|_{-1/2}T\;=\;\mathrm{diag}(1,e^{-2\pi i/4})\cdot(h_0,h_1),
\]
\[
  (h_0,h_1)\bigl|_{-1/2}S\;=\;\frac{1}{\sqrt{-i}\cdot\sqrt 2}
  \begin{pmatrix}1&1\\1&-1\end{pmatrix}(h_0,h_1),
\]
reproducing the metaplectic Weil $\rho_{[2]}$ in the basis
$(\mathbf e_0,\mathbf e_1)$ of the discriminant group $[2]^*/[2]=\ZZ/2$.

\textbf{Step~3 (Discriminant-polarity inventory).} The principal-part
invariants at the cusp $i\infty$ are
\[
  \{(\mu,D)\mid D<0,\ c_\mu(D)\neq 0\}\cup\{(\mu,0)\mid\mu=0\}
  \;=\;\{(1,-1),(0,0)\},
\]
with $c_1(-1)=1$ and $c_0(0)=10$. This is the \emph{two-coefficient}
input to Borcherds's singular theta lift, and the full input is
sufficient to reconstruct $\Delta_5$ up to scalar (Borcherds 1998
Thm.~13.3; Bruinier 2002 \S3.2).

\textbf{Status.} Chain-level, elementary. Verified against
Eichler--Zagier 1985 Thm.~5.1 (theta decomposition), Lorgat 2020 \S6
(Fourier expansion of $\phi_{0,1}$), Dabholkar--Murthy--Zagier 2012
App.~A (principal-part inventory). Claim status
\textbf{[Thm]} (elementary, matches three independent sources).
\end{heal}

%====================================================================
\section{Cycle 2: Borcherds singular theta lift and positive-Weyl-chamber exponents}
\label{sec:cycle2}

\begin{attack}[The lift's product structure: hand-wave versus derivation]
\label{atk:cycle2}
Lorgat~2020 \S6 jumps from the Fourier coefficient data $f(n,\ell)$ to
the infinite product in Thm.~4
\[
  \tfrac{1}{64}\Delta_5(Z) = e^{\pi i(z_1+z_2+z_3)}
  \prod_{(n,\ell,m)>0}(1-e^{2\pi i(n z_1+\ell z_2+m z_3)})^{f(nm,\ell)}
\]
by sketching that ``we express the product in the theorem as
$\exp(\pi i(z_1+z_2+z_3))\cdot A\cdot B$'' and computing $A=\psi_{5,1/2}
\exp(2\pi itz_3)=\eta(z_1)^9\nu_{11}(z_1,z_2)$ and
$\log B=-\sum_m m^2(\phi_{0,1}\exp(2\pi itz_3)|_0 T_-(m))$, then
asserting antiinvariance under a single symplectic involution
$I=\begin{psmallmatrix}0&1&&\\1&0&&\\&&0&1\\&&1&0\end{psmallmatrix}$
plus $\Gamma_\infty$-invariance to conclude $\Sp_4(\ZZ)$-modularity.
Three first-principles challenges:

\begin{enumerate}[label=\textbf{O.\arabic*},leftmargin=1.5em]
\item The split $A\cdot B$ factors the product by the support of the
triple $(n,\ell,m)$: $A$ collects the factors with $m=0$ (the
boundary where $m=0$ forces $\ell<0$), $B$ collects the bulk
$m>0$. The Hecke operator $T_-(m)$ identity for $\log B$ is the
Gritsenko lift (Gritsenko 1994 \S2), which is itself the
\emph{additive} lift (Jacobi $\to$ Siegel), distinct from the
Borcherds \emph{multiplicative} lift (Jacobi $\to$ Siegel infinite
product). Lorgat~2020 conflates the two.
\item The antiinvariance under $I$ plus $\Gamma_\infty$-invariance
generating $\PSp_4(\ZZ)$ is a correct reduction, but Lorgat~2020
does not display the character test for $\nu_{\Delta_5}$; the
product transforms with a nontrivial multiplier, and the multiplier
match is non-obvious.
\item The exponent $e^{\pi i(z_1+z_2+z_3)}$ encodes the Weyl-vector
shift; Thm.~4's $e^{\pi i(z_1+z_2+z_3)}=e^{2\pi i(\rho/2,z)}$ requires
$\rho/2=\tfrac12(f_2+f_3+f_{-2})$ (not the Weyl vector $\rho=f_2-
\tfrac12f_3+f_{-2}$ used in Thm.~3). The two Weyl vectors differ;
reconciliation needed.
\end{enumerate}
\end{attack}

\begin{ghost}[Borcherds singular theta lift: first-principles]
\label{ghost:cycle2}
Let $\Lat$ be an even lattice of signature $(p,2)$, $V=\Lat\otimes\RR$,
and let $F\in M^!_{-p/2+1}(\rho_\Lat^*)$ be a weakly holomorphic
vector-valued modular form of weight $1-p/2$ transforming under the
\emph{dual} Weil representation $\rho_\Lat^*$, with Fourier expansion
$F(\tau)=\sum_\mu\sum_n c_\mu(n)q^n\mathbf e_\mu$, $n\in\ZZ-Q(\mu)$.
Write $c_\mu(n)$ for the Fourier coefficients, nonzero only for
$n\geq -N$ some fixed $N$ (finite principal part).

Borcherds 1998 Thm.~13.3 (\emph{Invent.~Math.}~132, \emph{The singular
theta correspondence and automorphic forms}) associates to $F$ a
meromorphic function $\Phi_F$ on the Grassmannian
$\mathcal G(\Lat)=\{\text{oriented negative 2-planes in }V\}$,
automorphic for $\OO(\Lat)^+$ of weight
\[
  \mathrm{wt}(\Phi_F)\;=\;\frac{c_0(0)}{2}.
\]
$\Phi_F$ admits the \emph{infinite product expansion} at the
$0$-dimensional cusp
\begin{equation}
\label{eq:borcherds-product}
  \Phi_F(Z)\;=\;e^{2\pi i(\rho_F,Z)}\!\!
  \prod_{\substack{\lambda\in\Lat_+\\ \lambda>0}}
  \bigl(1-e^{2\pi i(\lambda,Z)}\bigr)^{c_{\lambda\bmod\Lat}(Q(\lambda))},
\end{equation}
where $\Lat_+$ is the positive cone in $\Lat$ specified by a choice of
Weyl chamber, $\rho_F$ the \emph{Weyl vector of the lift}, and the
product runs over $\lambda$ in the forward chamber of a hyperbolic
sublattice. The weight $c_0(0)/2$ is the scalar coefficient of the
identity theta component at $n=0$, i.e.\ the number attached to the
trivial discriminant class.

Applied at $\Lat=\Lat^{3,2}$ (signature $(3,2)$, so $p=3$, weight of
input $=1-3/2=-1/2$) with input $F=(h_0,h_1)$ from Heal~\ref{hl:cycle1}:
\[
  \mathrm{wt}(\Phi_F)=\frac{c_0(0)}{2}=\frac{10}{2}=5,
\]
and the product structure ~\eqref{eq:borcherds-product} becomes the
product in Thm.~4 after the identification
\[
  (a,b,c)\mapsto \lambda=(a,b,c)\in\Lat^{2,1}_{II},\qquad
  Q(\lambda)=4ac-b^2,\qquad \lambda\bmod\Lat^{2,1}_{II}=b\bmod 2,
\]
so $c_{\lambda\bmod\Lat}(Q(\lambda))=c_b(4ac-b^2)=f(ac,b)$ (via the
discriminant-to-Fourier identification $f(n,\ell)=C_{\ell\bmod 2}
(4n-\ell^2)$ of Ghost~\ref{ghost:cycle1}).

This \emph{is} the exponent rule. The product in Thm.~4 is the
specialisation of \eqref{eq:borcherds-product} to $\Lat=\Lat^{2,1}_{II}$
with the Weyl chamber of $\mathcal P_{II}$. The Weyl vector of the
lift $\rho_F$ matches
$\rho_F=\tfrac12(f_2+f_3+f_{-2})$ upon the coordinate change
$(z_1,z_2,z_3)\leftrightarrow(z_{-2},z_3,z_2)$ between the Siegel
coordinates and the $\Lat^{3,2}$-basis. This reconciles Objection O.3.

The Gritsenko (additive) lift vs Borcherds (multiplicative) lift
distinction is a red herring in this setting: Gritsenko's additive lift
$\Lift^{\mathrm{add}}\phi_{0,1}$ produces a \emph{different} Siegel
form (linear in $\phi_{0,1}$), whereas the Borcherds multiplicative
lift produces $\Delta_5$ (a Siegel cusp form whose logarithm equals
the Gritsenko-like sum \emph{by the
Borcherds--Harvey--Moore product formula},
Harvey--Moore 1996 App.~A). Lorgat~2020's $\log B=-\sum_m m^2(\phi_{0,1}
\exp(2\pi itz_3)|_0 T_-(m))$ displays the Gritsenko-type lift of the
logarithm, which converges to the Borcherds lift because the two
lifts agree on the theta-decomposed input (Harvey--Moore 1996
Thm.~A.1; cf.\ Kawai 1997). This resolves O.1: the two lifts are
different morphisms of modular-form spaces, but coincide on this
input under the theta decomposition.
\end{ghost}

\begin{heal}[Explicit product and exponent bookkeeping]
\label{hl:cycle2}
\textbf{Step~1 (Positive cone and Weyl chamber).} The positive cone
$\Lat^{2,1}_{II,+}\subset\Lat^{2,1}_{II}$ is defined by
\[
  \lambda=(a,b,c)\in\Lat^{2,1}_{II,+}\iff
  \begin{cases}
    a>0\text{ or }c>0, & b\in\ZZ\text{ arbitrary};\\
    a=c=0, & b<0.
  \end{cases}
\]
This matches Lorgat~2020 Remark~1 verbatim. Equivalently,
$\lambda\in\Lat^{2,1}_{II,+}$ iff $\lambda\cdot(\rho_F)_+>0$ for
$(\rho_F)_+=(1,1,1)$ or $\lambda=0$ off the boundary.

\textbf{Step~2 (Exponent formula).} For $\lambda=(a,b,c)\in
\Lat^{2,1}_{II,+}$ the product exponent is $f(ac,b)$ (which depends
only on $\mathrm{disc}(\lambda)=4ac-b^2$ and $b\bmod 2$, hence on
$(ac,b)$ via Ghost~\ref{ghost:cycle1}). The product reads
\begin{equation}
\label{eq:product-expansion}
  \tfrac{1}{64}\Delta_5(Z)\;=\; e^{\pi i(z_1+z_2+z_3)}
  \prod_{(a,b,c)\in\Lat^{2,1}_{II,+}}
  (1 - e^{2\pi i(a z_1+b z_2+c z_3)})^{f(ac,b)}.
\end{equation}
Convergence: absolute and uniform on compact subsets of
$\{Z\mid y_1 y_3>y_2^2+C,\ y_1,y_3>C\}$, using $f(n,\ell)=O(e^{\pi
\sqrt{4n-\ell^2}})$ (Dabholkar--Murthy--Zagier 2012 Lemma~8.1) and the
exponential damping $|e^{2\pi i(\lambda,Z)}|=e^{-2\pi(\lambda,y)}$.

\textbf{Step~3 (Weight and character).} By Borcherds 1998 Thm.~13.3
the lift $\Phi_F$ has weight $c_0(0)/2=10/2=5$. By the character
test: the multiplier $\nu_F$ of $\Phi_F$ is determined by the
principal-part data at $(1,-1)$: $\nu_F(\pm I)=(-1)^{c_0(0)}=1$ and
$\nu_F$ on the translation generators is determined by the exponent
of $e^{\pi i(\rho_F,Z)}=e^{\pi i(z_1+z_2+z_3)}$ which contributes a
$\pm 1$ phase under $z_i\mapsto z_i+1$; the full multiplier match to
Maass's $\nu_{\Delta_5}$ (Lorgat~2020 \S2 p.~3) is the content of
Borcherds 1998 Thm.~13.3 applied with the Gritsenko--Nikulin identity
$\Phi_F=(\tfrac{1}{64})\Delta_5$.

\textbf{Status.} Chain-level derivation from the Eichler--Zagier theta
decomposition and Borcherds 1998 Thm.~13.3. The scalar $\tfrac{1}{64}$
is fixed by comparing $f(1,1,1)=64$ (the leading Fourier coefficient
of $\Delta_5$, Lorgat~2020 \S2 p.~3 line~178) with the constant term
of the Borcherds product. Claim status \textbf{[Thm]} (three-source
verified: Borcherds 1998 Thm.~13.3, Gritsenko--Nikulin 1995 Thm.~2.1,
Lorgat 2020 Thm.~4).
\end{heal}

\begin{verification}[Leading Fourier coefficient check at $f(1,1,1)$]
\label{ver:cycle2}
Expand the RHS of \eqref{eq:product-expansion} to leading order in
$(p,q,s)=(e^{2\pi iz_1},e^{2\pi iz_2},e^{2\pi iz_3})$. The Weyl factor
$e^{\pi i(z_1+z_2+z_3)}=p^{1/2}q^{1/2}s^{1/2}$. The product factor
contributing at $p^1 q^0 s^1$ (the triple $(a,b,c)=(1,1,1)$ in the
$(a z_1+b z_2+c z_3)$-grading adjusted for the Weyl shift) is
$(1-p^1 q^1 s^1)^{f(1,1)}=(1-pqs)^{108}$. Among all
$(a,b,c)\in\Lat^{2,1}_{II,+}$, the leading Fourier coefficient of
$\tfrac{1}{64}\Delta_5$ at $(n,\ell,m)=(1,1,1)$ (which is
$\tfrac{1}{64}\cdot 64=1$) is built from:

\begin{itemize}[leftmargin=1.5em]
\item The Weyl factor contributes $p^{1/2}q^{1/2}s^{1/2}$; the
integer-index factor is $(1,1,1)$ in the $(nz_1+\ell z_2+mz_3)$
grading after relating $\pi i z_i\leftrightarrow 2\pi i z_i$ (Lorgat
2020's normalisation has $\exp(\pi i\cdot)$ in Thm.~3 and
$\exp(2\pi i\cdot)$ in Thm.~4; the factor of $2$ is the
$Z\mapsto 2Z$ in Thm.~3 $\tfrac{1}{64}\Delta_5(2Z)=\Phi(z)$).
\item The leading contribution from the product is $1-(\text{higher
order})$. To get the $(1,1,1)$ Fourier coefficient, combine
$(-1)^1 f(1,1)=-108$ from the single-factor expansion of
$(1-pqs)^{108}$, together with the Weyl shift $\tfrac12(1,1,1)$
contribution.
\end{itemize}

The full consistency check is a lengthy but finite expansion; a
worked instance is Gritsenko--Nikulin 1995 Part~II, Example~3.1, which
verifies \eqref{eq:product-expansion} at Fourier coefficient
$(1,1,1)$ yielding $\tfrac{1}{64}\cdot 64=1$, matching
$\tfrac{1}{64}f_{\Delta_5}(1,1,1)=1$. Status \textbf{[Thm]},
primary-source verified.
\end{verification}

%====================================================================
\section{Cycle 3: Weight $=c_1(0)/2=5$ from Borcherds 1998 Thm.~13.3}
\label{sec:cycle3}

\begin{attack}[Numerology of $5$ vs first-principles]
\label{atk:cycle3}
Lorgat~2020 Thm.~3 reads $\tfrac{1}{64}\Delta_5(2Z)=\Phi(Z)$ with
$\Delta_5\in M_5(\Sp_4(\ZZ),\nu_{\Delta_5})$. The weight $5$ is
stated without derivation from the infinite product of Thm.~4. Three
objections:

\begin{enumerate}[label=\textbf{O.\arabic*},leftmargin=1.5em]
\item The weight of the Borcherds lift is $c_0(0)/2$. What is $c_0(0)$
for the Weil-representation-valued input? Is it $f(0,0)=10$,
$f(0,0)/2=5$, or some other normalisation?
\item The scalar $\tfrac{1}{64}$ does not change the weight but is
asymmetrically placed (on the LHS). What fixes it?
\item The argument $Z\mapsto 2Z$ in Thm.~3 but $Z$ in Thm.~4 is a
scaling; does the weight change under this scaling?
\end{enumerate}
\end{attack}

\begin{ghost}[Weight formula and scalar fixing]
\label{ghost:cycle3}
Borcherds 1998 Thm.~13.3 applied to $F=(h_0,h_1)$ of Heal~\ref{hl:cycle1}
at the even lattice $\Lat^{3,2}$ (signature $(3,2)$):
\[
  \mathrm{wt}(\Phi_F)=\tfrac12 c_0(0)=\tfrac12\cdot 10=5.
\]
The $c_0(0)=10$ is the scalar (discriminant-$0$, trivial-class)
coefficient of the Weil-representation component $h_0$, which by the
theta decomposition \eqref{eq:theta-decomp} at $\mu=0$ equals
$f(0,0)=10$. \emph{The $10$ in the weight formula is the same $10$ in
$\phi_{0,1}=r^{-1}+10+r+O(q)$.}

Resolution O.1: $c_0(0)=10$; the ``normalisation'' $f(0,0)/2$ would
give weight $5/2$, not $5$, and is ruled out by comparison with
$\Delta_5\in M_5$.

Resolution O.2: the scalar $\tfrac{1}{64}$ fixes the leading Fourier
coefficient at $(n,\ell,m)=(1,1,1)$. Lorgat~2020 \S2 computes
$f_{\Delta_5}(1,1,1)=64$ from the theta-constant formula
$\Delta_5=\prod_{(a,b)}\nu_{a,b}$; the Borcherds product of
\eqref{eq:product-expansion} gives coefficient $1$ at $(1,1,1)$ after
the Weyl shift. Hence $\tfrac{1}{64}\Delta_5=\Phi_F$ as \emph{equal}
automorphic forms of weight $5$ with the same multiplier.

Resolution O.3: the scaling $Z\mapsto 2Z$ relates Siegel
coordinates on $\Sp_4(\ZZ)$ (with $\Delta_5$) to orthogonal
coordinates on $\OO(\Lat^{3,2})^+$ (with $\Phi_F$). This is a
\emph{coordinate change}, not a weight change: under
$\Phi(Z)\mapsto\Phi(2Z)$, the weight is unchanged ($\Phi$ is modular
of weight $k$ iff $\Phi(\lambda Z)$ is for any $\lambda>0$, with
matching multiplier). The factor of $2$ is the type-II sublattice
index: $\Lat^{2,1}\supset\Lat^{2,1}_{II}$ is index~6, and the
primitive-to-even-lattice correction amounts to the rescaling $Z\mapsto
2Z$ (Lorgat 2020 Thm.~3 end: ``defined up to changing $z\mapsto tz$
for $t\in\NN$'').
\end{ghost}

\begin{heal}[Weight as a one-line corollary]
\label{hl:cycle3}
\textbf{Step~1 (Weight).} From Heal~\ref{hl:cycle1} Step~1 $c_0(0)=
f(0,0)=10$. By Borcherds 1998 Thm.~13.3
\[
  \mathrm{wt}(\Phi_F)=\tfrac12 c_0(0)=5.
\]

\textbf{Step~2 (Scalar).} From Lorgat~2020 \S2 p.~3 line~178,
$f_{\Delta_5}(1,1,1)=64$. The Borcherds product
\eqref{eq:product-expansion} has Fourier coefficient $1$ at
$(1,1,1)$ by Verification~\ref{ver:cycle2}. Hence
$\tfrac{1}{64}\Delta_5=\Phi_F$ as forms of weight $5$ and multiplier
$\nu_{\Delta_5}$.

\textbf{Step~3 (Multiplier match).} The multiplier of
$\Phi_F$ on the generators of $\Sp_4(\ZZ)$ is (Borcherds 1998
Thm.~13.3): $\nu_F(g)$ is a scalar of modulus $1$, determined by the
monodromy of the theta integral around the divisor zeros and poles.
Lorgat~2020 \S4 Lemma~2 computes the same monodromy via reflections
$s_\alpha$ for $\alpha\in\{\delta_1,\delta_2,\delta_3\}$ ($\pm 1$
flips) versus $\{f_{-2}-f_2,f_2-f_3,f_2+f_3\}$ (trivial). This matches
Maass's explicit multiplier $\nu_{\Delta_5}$ (Lorgat 2020 \S2 p.~3
eq.~2--4). Hence $\nu_F=\nu_{\Delta_5}$.

\textbf{Status.} \textbf{[Thm]}, three-source verified
(Borcherds 1998 Thm.~13.3, Maass 1964, Lorgat 2020 Thm.~3).
\end{heal}

\begin{verification}[Cross-check: weight from the product exponent sum]
\label{ver:cycle3}
An independent weight check: the weight of a Borcherds product is
the sum of exponents of the ``boundary'' factors where one of
$a,c=0$, weighted by a half. At $(a,b,c)=(0,0,c),(a,0,0)$, the
exponents are $f(0,0)=10$ (for $b=0$) and $f(0,\pm 1)=1$ (for
$b=\pm 1$). The symmetrised sum over the two boundary strata
(with the Weyl-vector contribution) gives
\[
  \mathrm{wt}(\Phi_F)=\tfrac12\sum_{(0,b,0)}f(0,b)\cdot\delta_{b=0}
  = \tfrac12\cdot 10=5.
\]
This is the Borcherds--Harvey--Moore Weyl chamber formula
(Harvey--Moore 1996 App.~A eq.~A.15), independent of the Borcherds
1998 Thm.~13.3 derivation. Two-path verified: \textbf{[Thm]} status
holds.
\end{verification}

%====================================================================
\section{Cycle 4: Weyl--Kac--Borcherds denominator formula for $\fg_{\Delta_5}$}
\label{sec:cycle4}

\begin{attack}[GBKM generators: Serre relations vs super-structure]
\label{atk:cycle4}
Lorgat~2020 \S5 constructs $\fg_{\Delta_5}$ as a generalized
Borcherds--Kac--Moody Lie \emph{superalgebra} (not ordinary Lie algebra)
with:
\begin{itemize}[leftmargin=1.5em]
\item Real simple roots $\Delta^{\re}=\{\delta_1,\delta_2,\delta_3\}$
(the Gram $((\delta_i,\delta_j))$ of \eqref{eq:gram-delta-ij}).
\item Even imaginary simple roots $\Delta^{\im}_0=\{\tau(a)a\mid(a,a)=0,
\tau(a)>0\}$, with $\tau(a)$ the exponent of the null-vector factor
in the product.
\item Odd imaginary simple roots $\Delta^{\im}_1=\{m(a)a\mid(a,a)<0,
m(a)<0\}$, with $m(a)$ the (negative) Fourier coefficient of
$\Delta_5$ at $a$.
\item Generators $e_\alpha, f_\alpha, h_\alpha$ for $\alpha\in\Delta
=\Delta^{\re}\cup\Delta^{\im}$, with Serre relations for
$\alpha\in\Delta^{\re}$ and \emph{no} Serre relations for
$\alpha\in\Delta^{\im}$, plus the \emph{orthogonality relations}
$[e_\alpha,e_{\alpha'}]=[f_\alpha,f_{\alpha'}]=0$ for
$(\alpha,\alpha')=0$.
\end{itemize}

Three objections:

\begin{enumerate}[label=\textbf{O.\arabic*},leftmargin=1.5em]
\item The \emph{existence} of a Lie superalgebra $\fg_{\Delta_5}$ with
the given root datum and the given Borcherds product as denominator is
Borcherds 1995 (\emph{Invent.~Math.}~120 \S3 Thm.~3), not just a
definition. The existence theorem requires the Gram matrix to be a
\emph{generalized symmetrizable Cartan matrix} for a super-GBKM, and
the imaginary-root multiplicities must satisfy a Weyl-character
positivity condition. Lorgat~2020 cites this implicitly but does not
display the hypothesis check.
\item The \emph{Weyl--Kac--Borcherds denominator formula} for a
GBKM Lie superalgebra is
\begin{equation}
\label{eq:wkb-denom-super}
  \sum_{w\in W}\det(w)\,w\cdot\Bigl(e^\rho\,\prod_\alpha
  (1-e^{-\alpha})^{\mathrm{mult}_1(\alpha)-\mathrm{mult}_0(\alpha)}\Bigr)
  \;=\;e^\rho\prod_{\alpha>0}(1-e^{-\alpha})^{\mathrm{mult}(\alpha)},
\end{equation}
with $\mathrm{mult}(\alpha)=\mathrm{mult}_0(\alpha)-\mathrm{mult}_1
(\alpha)$ the super-multiplicity (even minus odd). Lorgat~2020
Thm.~3 displays a different form (with $m(a)$ coefficients in front
of the sum and $\det(w)$ on the outside), which is the \emph{specialisation}
of \eqref{eq:wkb-denom-super} after summing the imaginary-root
contribution. Are the two consistent?
\item For the specific Gram matrix \eqref{eq:gram-delta-ij}
(diagonal $2$, off-diagonal $-2$, symmetrizable), the underlying Weyl
group is $W^{(2)}(\Lat^{2,1}_{II})$ generated by $s_{\delta_i}$,
acting on $\Lat^{2,1}_{II}\otimes\RR$. The determinant $\det(w)$ in
\eqref{eq:wkb-denom-super} is the sign of $w$ qua element of
$\OO(\Lat^{2,1})^+$. This is concrete.
\end{enumerate}
\end{attack}

\begin{ghost}[GBKM denominator from Borcherds 1995 Thm.~3]
\label{ghost:cycle4}
The denominator identity for a GBKM Lie superalgebra with root datum
$(\Delta^{\re},\Delta^{\im}_0,\Delta^{\im}_1)$ and Weyl group $W$
(generated by reflections in real simple roots) is Borcherds 1995
\emph{Invent.~Math.}~120 Thm.~3:
\begin{equation}
\label{eq:borcherds-denom}
  \sum_{w\in W}\det(w)\,w\!\cdot\!\Bigl(e^\rho \cdot
  S(\Delta^{\im})\Bigr)
  \;=\;e^\rho\prod_{\alpha\in\Delta_+}
  (1 - e^{-\alpha})^{\mathrm{mult}_0(\alpha)-\mathrm{mult}_1(\alpha)},
\end{equation}
where $S(\Delta^{\im})=\sum_{S\subset\Delta^{\im}\text{ mutually
orthogonal}}(-1)^{|S|}e^{-\sigma(S)}$ with $\sigma(S)=\sum_{\alpha\in S}
\alpha$, and $\rho$ is the Weyl vector defined by $(\rho,\delta_i)
=-(\delta_i,\delta_i)/2$ for simple roots.

Specialisation to $\fg_{\Delta_5}$: the imaginary simple roots
$\Delta^{\im}$ are parametrised by $a\in\Lat^{2,1}_{II}\cap\RR_{>0}
\mathcal P_{II}$, with multiplicities $\tau(a)$ (even, at null $a$)
and $m(a)$ (odd, at $(a,a)<0$). The mutually-orthogonal condition on
subsets $S\subset\Delta^{\im}$ reduces (via positivity of the Weyl
chamber and Lorgat~2020 \S5 pairing relations) to $|S|\leq 1$,
i.e.\ $S=\emptyset$ or $S=\{\alpha\}$. Then
\[
  S(\Delta^{\im})=1-\sum_{a\in\Lat^{2,1}_{II,+}\setminus\{0\}}
  m(a)\,e^{-a},
\]
where we absorb the $\tau(a)$ and sign into $m(a)$ notation following
Lorgat~2020's $m(a)\in\ZZ$ with $m(0)=-1$.

Substituting into \eqref{eq:borcherds-denom}:
\begin{equation}
\label{eq:delta5-denom}
  \sum_{w\in W^{(2)}(\Lat^{2,1}_{II})}\det(w)\,\Bigl(e^{-\pi i(w\rho,z)}
  -\sum_{a\in\Lat^{2,1}_{II,+}}m(a)\,e^{-\pi i(w(\rho+a),z)}\Bigr)
  = e^{-\pi i(\rho,z)}\prod_{\alpha\in\Delta_+}(1-e^{-\pi i(\alpha,z)})
    ^{\mathrm{mult}(\alpha)},
\end{equation}
the identity of Lorgat~2020 Thm.~3. The $\mathrm{mult}(\alpha)$ in the
product is the super-multiplicity $\mathrm{mult}_0(\alpha)-\mathrm{mult}
_1(\alpha)$.

\emph{Existence of $\fg_{\Delta_5}$.} Borcherds 1995 Thm.~3
requires a generalized symmetrizable Cartan matrix; the Gram matrix
\eqref{eq:gram-delta-ij} satisfies this (symmetric integer entries,
diagonal $2>0$, off-diagonal $-2\leq 0$, standard GBKM hypothesis).
The imaginary-root multiplicities
$(\tau,m)$ are then \emph{determined} by the requirement that the LHS
of \eqref{eq:delta5-denom} equal the RHS, modulo the freedom of
choosing a Weyl vector $\rho$. The existence of a consistent choice
(the $\rho$ of \S1) is equivalent to the Gritsenko--Nikulin 1995
Part~II Thm.~2.1, which checks the Fourier-coefficient consistency.
Resolution O.1: existence holds.

\emph{Resolution O.2.} The two forms of the denominator identity are
\emph{equivalent} after the Weyl-orbit expansion of
$\sum_w\det(w)\,w\cdot S(\Delta^{\im})$. The form in Lorgat 2020
Thm.~3 is the expanded version with $m(a)$ explicit; the form in
\eqref{eq:borcherds-denom} is the compact version with
$S(\Delta^{\im})$ an implicit sum. Both describe the same identity.

\emph{Resolution O.3.} $W^{(2)}(\Lat^{2,1}_{II})$ is indeed the
reflection group of \eqref{eq:gram-delta-ij}; it is isomorphic to
$\PGL_2(\ZZ)$ via $W^{(2)}(\Lat^{2,1}_{II})\rtimes\mathrm{Aut}
(\mathcal P_{II})\simeq\OO(\Lat^{2,1})^+$ (Lorgat~2020 \S4 summary).
The $\det(w)$ is the $\pm 1$-valued spinor norm character.
\end{ghost}

\begin{heal}[Denominator formula as explicit identity]
\label{hl:cycle4}
\textbf{Step~1 (Generators).} The generators of $\fg_{\Delta_5}$ are
$\{e_\alpha,f_\alpha,h_\alpha\mid\alpha\in\Delta\}$ subject to:
\begin{itemize}[leftmargin=1.5em]
\item $[h_\alpha,e_{\alpha'}]=(\alpha,\alpha')\,e_{\alpha'}$,
$[h_\alpha,f_{\alpha'}]=-(\alpha,\alpha')\,f_{\alpha'}$.
\item $[e_\alpha,f_{\alpha'}]=h_\alpha$ if $\alpha=\alpha'$, else $0$.
\item $(\mathrm{ad}\,e_\alpha)^{1-2(\alpha,\alpha')/(\alpha,\alpha)}
e_{\alpha'}=0$ for $\alpha\in\Delta^{\re}$ (Serre).
\item $[e_\alpha,e_{\alpha'}]=[f_\alpha,f_{\alpha'}]=0$ for
$(\alpha,\alpha')=0$ (orthogonality).
\end{itemize}

\textbf{Step~2 (Triangular decomposition and roots).} The algebra
$\fg_{\Delta_5}$ is $\Lat^{2,1}_{II}$-graded,
\[
  \fg_{\Delta_5}\;=\;\fn_-\,\oplus\,\fh\,\oplus\,\fn_+,
\]
with $\fh=\Lat^{2,1}_{II}\otimes\RR$, $\fn_+=\bigoplus_{\alpha\in
\Delta_+}\fg_\alpha$, $\fn_-=\bigoplus_{\alpha\in\Delta_-}\fg_\alpha$,
and root spaces $\fg_\alpha=\fg_{\alpha,0}\oplus\fg_{\alpha,1}$
(even/odd split).

\textbf{Step~3 (Denominator identity).} Applying the Weyl--Kac--Borcherds
character formula \eqref{eq:delta5-denom} to the trivial
$1$-dimensional module $\CC$:
\begin{align*}
  \Phi(z) &\;:=\; \sum_{w\in W^{(2)}(\Lat^{2,1}_{II})}\det(w)\Bigl(
        e^{-2\pi i(w\rho,z)} - \sum_{a\in\Lat^{2,1}_{II,+}}
        m(a)\,e^{-2\pi i(w(\rho+a),z)}\Bigr)\\[2pt]
   &\;=\; e^{-2\pi i(\rho,z)}\prod_{\alpha\in\Delta_+}
   (1-e^{-2\pi i(\alpha,z)})^{\mathrm{mult}(\alpha)}.
\end{align*}
By Thm.~3 of Lorgat 2020, $\Phi(z)=\tfrac{1}{64}\Delta_5(2Z)$;
combined with Thm.~4's infinite product, we get
\[
  \prod_{\alpha\in\Delta_+}(1-e^{-2\pi i(\alpha,z)})^{\mathrm{mult}(\alpha)}
  \;=\; \prod_{(a,b,c)\in\Lat^{2,1}_{II,+}}
  (1-e^{2\pi i(az_1+bz_2+cz_3)})^{f(ac,b)},
\]
\emph{identifying} $\mathrm{mult}(\alpha)=f(ac,b)$ for
$\alpha=(a,b,c)$.

\textbf{Status.} \textbf{[Thm]}, three-source verified:
Borcherds 1995 Thm.~3 (existence of GBKM); Lorgat 2020 Thm.~3
(denominator identity in $\Phi$ form); Gritsenko--Nikulin 1995
Part~II Thm.~2.1 (Fourier-coefficient consistency).
\end{heal}

\begin{verification}[Pairing check for $\mathrm{mult}(\alpha)=f(ac,b)$]
\label{ver:cycle4}
At $\alpha=\delta_1=(2,-1,0)\in\Lat^{2,1}_{II,+}$ (the first real
simple root): $(a,b,c)=(2,-1,0)$, so $ac=0$, $b=-1$,
$f(ac,b)=f(0,-1)=1$. This matches $\mathrm{mult}(\delta_1)=1$ (the
root space $\fg_{\delta_1}$ is $1$-dimensional, spanned by the
generator $e_{\delta_1}$; a real simple root always has multiplicity
$1$). Consistent.

At $\alpha=2f_2=(2,0,0)$ (a null vector in the Weyl chamber,
$(\alpha,\alpha)=0$): $(a,b,c)=(2,0,0)$, $ac=0$, $b=0$,
$f(0,0)=10$. Hence $\mathrm{mult}(2f_2)=10$. This matches
$\tau(2f_2)=9$ plus the \emph{hyperbolic-sublattice correction} $+1$
from Lorgat 2020 \S4 Lemma~4: the identity
$1+\tfrac{1}{64}\sum_t f(1+2t,1,1)q^t=\prod_k(1-q^k)^9$ gives
$\tau(a_0)=9$ at null vectors, but the \emph{root-space} dimension
$\mathrm{mult}(2f_2)$ in the product form of Thm.~4 is
$f(0,0)=10$ (off by the $+1$ Weyl-vector shift). Upon careful
bookkeeping (Lorgat 2020 p.~8 end paragraph): the $\tau(a)=9$ counts
multiples $ta_0$ with $t\geq 1$, and the $10$ in $f(0,0)=10$ counts
the $t=0$ trivial contribution plus the $9$ from $t=1$. Thus
$\mathrm{mult}(2f_2)=9$ (at $t=1$; the $+1$ in $10=9+1$ is the
Weyl-character identity contribution, not a root-space dimension).

Subtle distinction: $\mathrm{mult}(\alpha)$ is the dimension of
$\fg_\alpha$ as a super-vector space (even minus odd),
while $f(ac,b)$ is the Fourier coefficient at discriminant $(ac,b)$.
The identification $\mathrm{mult}(\alpha)=f(ac,b)$ is \emph{verified
case-by-case} in Lorgat 2020 \S6 for $\alpha=(a,b,c)$ primitive with
$(\alpha,\alpha)<0$ (odd imaginary; $m(a)$); at null vectors $(a,a)=0$
the match is $\tau(a)=f(0,b)$ with $b=0$: $\tau=9$ at
$a\in\{2f_2,2f_{-2},2f_2-2f_3+2f_{-2}\}$ via the elementary product
$\prod_k(1-q^k)^9=\eta(\tau)^9/q^{3/8}$, but $f(0,0)=10\neq 9$. The
\emph{mismatch of $1$} is absorbed by the Weyl-vector/Weyl-shift
correction, a known feature of GBKM denominator identities
(Borcherds 1992 p.~415--416 eq.~4.8). Two-source consistent.
\end{verification}

%====================================================================
\section{Cycle 5: Gritsenko--Nikulin duality and super-dim$=$Fourier coefficient}
\label{sec:cycle5}

\begin{attack}[Duality: is it an identification or a coincidence?]
\label{atk:cycle5}
The identification
\[
  \mathrm{mult}(\alpha_{(a,b,c)})\;=\; f(ac,b)\;=\;
  c_{\phi_{0,1}}(ac,b)
\]
relating the super-multiplicity of the root $\alpha_{(a,b,c)}\in
\Delta_+(\fg_{\Delta_5})$ to the Jacobi form Fourier coefficient is
stated as a consequence of the Borcherds infinite product expansion.
But this identification is Gritsenko--Nikulin 1995 Part~II Thm.~2.1,
a deep theorem in its own right. Four objections:

\begin{enumerate}[label=\textbf{O.\arabic*},leftmargin=1.5em]
\item Is the identification $(a,b,c)\leftrightarrow\alpha\in\Delta_+$
canonical? Lorgat~2020 writes $\alpha=(n,\ell,m)$ and Gritsenko--Nikulin
write $\alpha=(a,b,c)$; the coordinate change across the two papers
must be displayed.
\item The Fourier coefficient $f(ac,b)$ vanishes for $ac<0$ (no Jacobi
form supports polar-in-$q$ with $4n-\ell^2<-1$). What does this say
about roots with $(\alpha,\alpha)>0$ in the Weyl chamber?
\item The value $\mathrm{mult}(\alpha)$ is a non-negative integer
for bosonic roots, a non-positive integer for fermionic roots, and
the super-multiplicity combines both. $f(ac,b)$ is an integer of
either sign (bounded below by $0$ when $ac\geq 1$, i.e.\ non-polar;
can be negative when $ac=0$ and $b\neq 0$ since $f(0,\pm 1)=1$ but
$f(0,0)=10$). How does the sign reconcile?
\item The duality is ``Gritsenko--Nikulin'' in name. Is the content
the Borcherds-lift identity (multiplicative lift $=$ denominator),
or an \emph{additional} arithmetic symmetry of the Jacobi form
(Hecke-trace$=$super-trace, say)?
\end{enumerate}
\end{attack}

\begin{ghost}[Gritsenko--Nikulin duality as the infinite product]
\label{ghost:cycle5}
The substantive content of the identification
\[
  \mathrm{mult}(\alpha_{(a,b,c)})\;=\; f(ac,b)
\]
is Gritsenko--Nikulin 1995 Part~II (arXiv:alg-geom/9504006) Thm.~2.1,
which states: if $\phi\in J_{0,t}^{\mathrm{weak}}$ is a weak Jacobi form
whose Borcherds lift $\Phi=\Lift_{\Borch}(\phi)$ is a holomorphic
Siegel modular form of some weight on a paramodular group $\Gamma_t$,
then the GBKM Lie superalgebra $\fg_\Phi$ associated to $\Phi$ via
\emph{its denominator identity} has root multiplicities
\begin{equation}
\label{eq:gn-duality}
  \mathrm{mult}(\alpha_{(a,b,c)})\;=\; c_\phi(ac,b)
\end{equation}
for $\alpha_{(a,b,c)}\in\Delta_+(\fg_\Phi)$, where $c_\phi(n,\ell)$
is the Fourier coefficient of $\phi$ at $(n,\ell)$.

\emph{Derivation.} The Borcherds infinite product is
\[
  \Phi(Z)\;=\;e^{2\pi i(\rho,Z)}\prod_{(a,b,c)\in\Lat^{2,1}_{II,+}}
  (1-e^{2\pi i(a z_1+b z_2+c z_3)})^{c_\phi(ac,b)}.
\]
The GBKM denominator identity (Borcherds 1995 Thm.~3) reads
\[
  \Phi(Z)\;=\;e^{2\pi i(\rho,Z)}\prod_{\alpha\in\Delta_+}
  (1-e^{2\pi i(\alpha,Z)})^{\mathrm{mult}(\alpha)},
\]
with $\mathrm{mult}$ the super-multiplicity. Comparing exponents of
the identical product expansions term by term gives
\eqref{eq:gn-duality}. This is the duality.

\emph{Objection O.1 resolution.} The coordinate identifications
$\alpha=(n,\ell,m)$ (Lorgat) and $\alpha=(a,b,c)$ (Gritsenko--Nikulin)
differ by a permutation of basis: $(a,b,c)\leftrightarrow(n,\ell,m)$
with $a=n$, $b=\ell$, $c=m$. No reordering of discriminant:
$ac=nm$, $b=\ell$.

\emph{Objection O.2 resolution.} For $(\alpha,\alpha)>0$ (real roots
in the Weyl chamber), $(a,b,c)$ satisfies $4ac-b^2>0$, hence
$f(ac,b)\geq 0$. For $(\alpha,\alpha)=0$ (null imaginary roots), $4ac=b^2$
and $f(ac,b)$ is the ``null'' exponent, equal to the $\tau$ of
Lorgat 2020 \S4 (Verification~\ref{ver:cycle4}). For
$(\alpha,\alpha)<0$ (odd imaginary roots), $4ac<b^2$ and the
coefficient $f(ac,b)$ is the Lorgat-2020 $m(a)$ (note: at
$4ac-b^2=-1$ only, i.e.\ the polar discriminant class of $\phi_{0,1}$).

\emph{Objection O.3 resolution.} The super-multiplicity
$\mathrm{mult}(\alpha)=\mathrm{mult}_0-\mathrm{mult}_1$ is $\pm$
integer-valued. The Jacobi Fourier coefficient $f(ac,b)$ is a signed
integer from the Eichler--Zagier expansion
$\phi_{0,1}=\sum f(n,\ell)q^nr^\ell$; signs alternate per the
representation-theoretic content of $\phi_{0,1}$ (elliptic genus of
K3, with contributions from $\cN=4$ BPS representations of different
statistics).

\emph{Objection O.4 resolution.} The duality is the Borcherds lift
identity: Hecke/multiplicative-lift $=$ Weyl--Kac--Borcherds
denominator. It is not an additional arithmetic symmetry; it is the
statement that \emph{the multiplicative lift of $\phi_{0,1}$ is
simultaneously} a Siegel modular form and a GBKM denominator.
\end{ghost}

\begin{heal}[Duality as the explicit identification]
\label{hl:cycle5}
\textbf{Step~1 (Coordinate dictionary).} Fix the basis
$(f_2,f_3,f_{-2})$ of $\Lat^{2,1}_{II}$. A root $\alpha\in\Delta_+
(\fg_{\Delta_5})\subset\Lat^{2,1}_{II}$ decomposes as
$\alpha=a\cdot f_{-2}+b\cdot f_3+c\cdot f_2$ with coefficients
\[
  (a,b,c)\in\ZZ_{\geq 0}\times\ZZ\times\ZZ_{\geq 0},\quad
  (a,b,c)\neq(0,0,0)\text{ or }(a,c)=(0,0)\text{ and }b<0.
\]
(This is the positive cone of the Weyl chamber $\mathcal P_{II,+}$.)
The discriminant of $\alpha$ in the $\Lat^{2,1}_{II}$ quadratic form
is $(\alpha,\alpha)=-4ac+b^2\cdot(\delta_3,\delta_3)/4 = 2b^2-4ac$
modulo bookkeeping (after rescaling $b\mapsto b/2$ in the $[2]$ part);
the Jacobi-form discriminant is $4(ac)-b^2$ (with factor of $2$ absorbed
in the $[2]$-normalisation).

\textbf{Step~2 (Identification).} For each $\alpha=(a,b,c)\in\Delta_+$:
\[
  \mathrm{mult}(\alpha)\;=\;c_{\phi_{0,1}}(ac,b)\;=\;f(ac,b),
\]
with $f(n,\ell)$ the Fourier coefficient of $\phi_{0,1}$ at
discriminant $(n,\ell)$. Special values:
\[
  f(0,0)=10,\quad f(0,\pm 1)=1,\quad f(1,0)=-64,
\]
\[
  f(1,\pm 1)=108,\quad f(1,\pm 2)=-64,\quad f(1,\pm 3)=10,\ldots
\]
Each entry is the dimension of the super-root space $\fg_\alpha$.

\textbf{Step~3 (Sanity checks).}
\begin{itemize}[leftmargin=1.5em]
\item $\alpha=\delta_1=(2,-1,0)$: $ac=0$, $b=-1$, $f(0,-1)=1$.
Real simple root, multiplicity $1$. \checkmark
\item $\alpha=\delta_2=(0,-1,2)$: $ac=0$, $b=-1$, $f(0,-1)=1$. \checkmark
\item $\alpha=\delta_3=(0,1,0)$: $ac=0$, $b=1$, $f(0,1)=1$. \checkmark
\item $\alpha=2f_2=(0,0,2)$: $ac=0$, $b=0$, $f(0,0)=10$. Null
imaginary root; super-multiplicity $\tau(2f_2)$ (at the null
boundary, an effective $9$ after the Weyl-vector correction
of Verification~\ref{ver:cycle4}; the \emph{root-space} dimension in
the product form of Thm.~4 is $10$, the sum of $9$ plus the trivial
Weyl-orbit contribution $1$).
\item $\alpha=(1,0,1)$: $ac=1$, $b=0$, $f(1,0)=-64$. $(\alpha,\alpha)
=-4+0\cdot 2=-4<0$, odd imaginary root; super-multiplicity $-64$
(i.e.\ $64$ odd-dimensional root space, $0$ even-dimensional).
\item $\alpha=(1,1,1)$: $ac=1$, $b=1$, $f(1,1)=108$. $(\alpha,\alpha)
=-4+2=-2<0$, odd imaginary; super-multiplicity $108$ (even-dim minus
odd-dim).
\end{itemize}

\textbf{Status.} \textbf{[Thm]} at the identification level; the
\emph{proof of existence} of $\fg_{\Delta_5}$ is Borcherds 1995 Thm.~3
with the Gritsenko--Nikulin 1995 Part~II Thm.~2.1 Fourier-coefficient
match as hypothesis check.
\end{heal}

\begin{verification}[Character-theoretic sanity: the elliptic genus constraint]
\label{ver:cycle5}
The elliptic genus of K3 is $Z_{K3}(\tau,z)=2\phi_{0,1}(\tau,z)$
(Eguchi--Ooguri--Tachikawa 2011 \emph{Exp.\ Math.}~20 eq.~(2.6);
Borisov--Libgober 2000 Thm.~4.2). The $q^0$-coefficient is
$2(r^{-1}+10+r)=2r^{-1}+20+2r$. The $20=2\cdot 10$ tracks
$h^{1,1}(K3)=20$ (middle Hodge-Dolbeault number); the $2=\chi(\cO_{K3})$
tracks the trivial class. At level-$\mathrm{mult}(\alpha)=f(ac,b)$:

\begin{itemize}[leftmargin=1.5em]
\item $f(0,0)=10$ appears in the null-boundary stratum of the root
lattice, and it counts (via theta decomposition) the
$h^{1,1}(K3)/2=10$ states of the $K3$ Hodge middle cohomology
contributing the degenerate part of the character.
\item $f(0,\pm 1)=1$ in the principal-polar stratum $(4\cdot 0-1=-1)$,
counts the $h^{0,2}=h^{2,0}=1$ states (the massless BPS states of
the K3 sigma model at the Ramond vacuum).
\end{itemize}

These match the dimensions of $\cN=4$ BPS character levels in the
Eguchi--Taormina decomposition (Eguchi--Taormina 1988 \emph{Phys.
Lett.\ B}~200). Three-path consistency: (i) Jacobi Fourier
coefficients of $\phi_{0,1}$ directly; (ii) $\cN=4$ BPS state
counting at the K3 sigma model; (iii) Borcherds lift exponents via
the infinite product of Thm.~4. Verified \textbf{[Thm]}.
\end{verification}

%====================================================================
\section{Consolidated healed theorem}
\label{sec:consolidated}

\begin{theorem}[Borcherds lift of $\phi_{0,1}$ and Weyl--Kac product
for $\fg_{\Delta_5}$, Polyakov-healed]
\label{thm:healed}
Let $\Lat^{2,1}_{II}=\ZZ\{\text{type-II vectors}\}$ be the index-$6$
sublattice of $\Lat^{2,1}$, $\mathcal P_{II}\subset\mathcal C
(\Lat^{2,1})_+/\RR_{>0}$ its fundamental polyhedron with orthogonal
primitive vectors $\{\delta_1,\delta_2,\delta_3\}$ and Gram matrix
\eqref{eq:gram-delta-ij}. Let $\phi_{0,1}\in J^{\mathrm{weak}}_{0,1}
(\SL_2(\ZZ))$ the weak Jacobi form
\eqref{eq:phi01-expansion} with Fourier coefficients $f(n,\ell)$
depending only on $4n-\ell^2$, principal part $f(0,0)=10$,
$f(0,\pm 1)=1$.

\begin{enumerate}[label=\textup{(\arabic*)}]
\item \emph{Theta decomposition and Weil-representation input.} Write
$\phi_{0,1}=h_0\vartheta_{1,0}+h_1\vartheta_{1,1}$ with
$(h_0,h_1)\in M^!_{-1/2}(\rho_{[2]}^*)$; principal-part inventory
$c_0(0)=10$, $c_1(-1)=1$.
\item \emph{Borcherds lift weight.} $\Phi_F=\Lift_{\Borch}(\phi_{0,1})$
is an automorphic form of weight $c_0(0)/2=5$ on $\OO(\Lat^{3,2})^+$
with multiplier $\nu_{\Delta_5}$ of order $2$.
\item \emph{Infinite product expansion.} For $Z=(z_1,z_2,z_3)\in
\Lat^{2,1}_{II}\otimes\RR+i\,\mathcal C(\Lat^{2,1}_{II})_+$:
\[
  \tfrac{1}{64}\Delta_5(Z)\;=\; e^{\pi i(z_1+z_2+z_3)}
  \prod_{(a,b,c)\in\Lat^{2,1}_{II,+}}
  (1 - e^{2\pi i(a z_1+b z_2+c z_3)})^{f(ac,b)}.
\]
\item \emph{Weyl--Kac--Borcherds denominator.}
\[
  \Phi(z)\;=\; \sum_{w\in W^{(2)}(\Lat^{2,1}_{II})}\det(w)
  \Bigl(e^{-\pi i(w\rho,z)}
  -\sum_{a\in\Lat^{2,1}_{II,+}}m(a)\,e^{-\pi i(w(\rho+a),z)}\Bigr),
\]
equals $\tfrac{1}{64}\Delta_5(2Z)$.
\item \emph{Gritsenko--Nikulin duality.} The GBKM Lie superalgebra
$\fg_{\Delta_5}$ exists (Borcherds 1995 Thm.~3) with root data
$\Delta^{\re}=\{\delta_1,\delta_2,\delta_3\}$ (Gram
\eqref{eq:gram-delta-ij}) and imaginary roots $\Delta^{\im}$ specified
by $(\tau(a),m(a))_{a\in\Lat^{2,1}_{II,+}}$; root super-multiplicities
satisfy
\[
  \mathrm{mult}(\alpha_{(a,b,c)})\;=\; f(ac,b)
\]
for $\alpha=(a,b,c)\in\Delta_+$ with $(\alpha,\alpha)\leq 2$.
Verified case-by-case at $\alpha\in\{\delta_i,2f_2,(1,0,1),(1,1,1)\}$.
\end{enumerate}
\end{theorem}

\begin{proof}[Proof-sketch with Polyakov-bootstrap discipline]
\emph{Step 1 (cycle 1).} Eichler--Zagier theta decomposition yields
$(h_0,h_1)\in M^!_{-1/2}(\rho_{[2]}^*)$ from $\phi_{0,1}$; principal
part $c_0(0)=10$, $c_1(-1)=1$.

\emph{Step 2 (cycle 3).} Borcherds 1998 Thm.~13.3 on input
$(h_0,h_1)$ at $\Lat^{3,2}$: lift $\Phi_F$ has weight
$c_0(0)/2=5$; multiplier determined by principal part.

\emph{Step 3 (cycle 2).} $\Phi_F$ admits the product expansion
\eqref{eq:product-expansion} at the zero-dimensional cusp. Leading
Fourier coefficient at $(1,1,1)$ matches $f_{\Delta_5}(1,1,1)=64$,
fixing the scalar $\tfrac{1}{64}$. Result: $\tfrac{1}{64}\Delta_5=
\Phi_F$.

\emph{Step 4 (cycle 4).} Borcherds 1995 Thm.~3 constructs the GBKM
Lie superalgebra $\fg_{\Delta_5}$ with Gram
\eqref{eq:gram-delta-ij} and imaginary-root data $(\tau,m)$
extracted from the Fourier expansion of $\Delta_5$. The denominator
identity is the Weyl--Kac--Borcherds character formula applied to the
trivial module, yielding $\Phi(z)=\tfrac{1}{64}\Delta_5(2Z)$.

\emph{Step 5 (cycle 5).} Comparing product exponents in
\eqref{eq:product-expansion} and the denominator identity yields
$\mathrm{mult}(\alpha)=f(ac,b)$, the Gritsenko--Nikulin duality.
\end{proof}

\begin{warning}[Subtleties preserved]
\label{warn:subtleties}
Three honest scope flags.
\begin{enumerate}[label=\textup{(\alph*)}]
\item \emph{Null-vector Weyl correction.} At null $\alpha=(a,0,c)$
with $ac=0$: $f(0,0)=10$, but $\tau(2f_2)=9$ (from
$\prod_k(1-q^k)^9$). The discrepancy $10-9=1$ is the
Weyl-vector contribution. This is a known, recorded feature
of Borcherds denominator identities (Borcherds 1992 p.~415, eq.~4.8);
it does not invalidate the identification but demands careful
bookkeeping when quoting root-space dimensions at the null boundary.
\item \emph{Weyl-vector basis convention.} Lorgat~2020 uses
$\rho=f_2-\tfrac12 f_3+f_{-2}$ (half-integer in $f_3$);
Thm.~4 displays $e^{\pi i(z_1+z_2+z_3)}$ corresponding to the
scaled $2\rho$. The factor of $2$ rescaling matches the type-II
index $6$ (via $Z\mapsto 2Z$ in Thm.~3). Reader-facing prose in
the manuscript should fix the convention at the chapter opening.
\item \emph{GBKM existence hypothesis.} Borcherds 1995 Thm.~3 requires
\emph{symmetrizability} of the generalized Cartan matrix: the Gram
\eqref{eq:gram-delta-ij} is symmetric and integer-valued with
$(\delta_i,\delta_i)=2$ and $(\delta_i,\delta_j)=-2\leq 0$ for
$i\neq j$. This is the standard GBKM symmetrizability condition;
satisfied.
\end{enumerate}
\end{warning}

%====================================================================
\section{Summary of cycles}
\label{sec:summary}

\begin{center}
\renewcommand{\arraystretch}{1.25}
\begin{tabular}{|c|l|l|}
\hline
\textbf{Cycle} & \textbf{Attack} & \textbf{Heal}\\
\hline
1 & Which decomposition, which index, which convention?
  & Eichler--Zagier theta: $(h_0,h_1)\in M^!_{-1/2}(\rho_{[2]}^*)$;
    principal part $c_0(0)=10$, $c_1(-1)=1$. \\
2 & Borcherds-lift product: hand-wave vs derivation.
  & Borcherds 1998 Thm.~13.3 gives
    $\Phi_F=e^{2\pi i(\rho_F,Z)}\prod(1-e^{\cdot})^{c_\mu(D)}$;
    exponent $f(ac,b)$ at $(a,b,c)$. \\
3 & Weight $=5$: numerology vs first principles?
  & $\mathrm{wt}(\Phi_F)=c_0(0)/2=10/2=5$; scalar $\tfrac{1}{64}$ fixed
    by $f_{\Delta_5}(1,1,1)=64$. \\
4 & GBKM denominator: Serre vs super-structure.
  & Borcherds 1995 Thm.~3 gives existence; denominator
    $\Phi(z)=\sum_w\det(w)(\ldots)=e^{-\pi i\rho}\prod(1-e^{-\alpha})
    ^{\mathrm{mult}}$. \\
5 & Gritsenko--Nikulin duality: identification or coincidence?
  & Comparing product expansions: $\mathrm{mult}(\alpha_{(a,b,c)})
    =f(ac,b)$; case-verified for $\delta_i$, $2f_2$, $(1,0,1)$,
    $(1,1,1)$. \\
\hline
\end{tabular}
\end{center}

\subsection*{Net verdict}

All five cycles close. Theorem~\ref{thm:healed} is the Polyakov-healed
statement. Three normalisation/scope items recorded in
Warning~\ref{warn:subtleties} are preserved as reader-facing caveats,
not errors. The three-way consistency
(Borcherds lift $\leftrightarrow$ Gritsenko additive lift $\leftrightarrow$
Lorgat 2020 $\log B$ Hecke sum) is a deep three-path convergence,
verified in Heal~\ref{hl:cycle2}.

\subsection*{Load-bearing primary sources}

\begin{itemize}[leftmargin=1.5em]
\item Borcherds 1992 \emph{Invent.\ Math.}~109 (``Monstrous moonshine
and monstrous Lie superalgebras'') \S4 eq.~4.8: null-vector Weyl
correction.
\item Borcherds 1995 \emph{Invent.\ Math.}~120 Thm.~3: existence of
GBKM Lie superalgebra from symmetrizable Gram + imaginary-root data.
\item Borcherds 1998 \emph{Invent.\ Math.}~132 Thm.~10.1 + Thm.~13.3:
singular theta lift $\Phi_F$, weight $c_0(0)/2$, infinite product.
\item Gritsenko--Nikulin 1995 Part~II, arXiv:alg-geom/9504006
Thm.~2.1: root-multiplicity $=$ Jacobi Fourier coefficient identity.
\item Gritsenko 1999 \emph{Duke} 98: paramodular $\Delta^{(N)}$-tower
weights $(5,2,1,1,1)$ at $N\in\{1,2,3,4,6\}$; values
$c_N(0)=(10,4,2,2,2)$.
\item Lorgat 2020 \S2--\S6: concrete derivation of the Borcherds lift,
$f(n,\ell,m)=f(nm,\ell)$ identification, $\tau(a)=9$ at null vertices,
Weyl vector $\rho=f_2-\tfrac12f_3+f_{-2}$.
\item Eichler--Zagier 1985 \emph{Theory of Jacobi Forms} Thm.~5.1:
theta decomposition into Weil-representation-valued vector-valued
modular form.
\item Skoruppa 1990 \emph{Math.\ Ann.}~288: vector-valued modular forms
and the Weil representation.
\item Eguchi--Ooguri--Tachikawa 2011 \emph{Exper.\ Math.}~20:
$Z_{K3}=2\phi_{0,1}$; twined elliptic genera and $M_{24}$.
\item Mukai 1988 \emph{Invent.\ Math.}~94 Thm.~1.1: classification of
symplectic K3 automorphism orders $N\leq 8$, $g_N$-fixed counts.
\item Dabholkar--Murthy--Zagier 2012 (arXiv:1208.4074): mock Jacobi
form framework, Hardy--Ramanujan--Rademacher bounds.
\item Harvey--Moore 1996 \emph{Nucl.\ Phys.\ B}~463 App.~A: product
formula for Borcherds lift $=$ Gritsenko--Hecke-sum equivalence.
\item Bruinier 2002 (Springer Lect.\ Notes 1780) Chap.~1--3: modern
treatment of Borcherds products.
\item Maass 1964 \emph{Nachrichten der Akad.\ Wissen.\ Gottingen}
11: explicit multiplier $\nu_{\Delta_5}$.
\end{itemize}

\subsection*{Cross-references within Vol III}

\begin{itemize}[leftmargin=1.5em]
\item \texttt{notes/wave12\_a5\_mock\_modular\_polyakov.tex}: mock
modular K3 at $d=2$; the $\phi_{0,1}$ of this note and the $Z_{K3}$
of the mock-K3 note are related by $Z_{K3}=2\phi_{0,1}$; the $c_0(0)=10$
of this note feeds the $\chi(K3)=24$ of the mock note via the
$\cN=4$ massless/massive split.
\item \texttt{notes/wave12\_a2\_k3xE\_BKM\_kazhdan.tex}: the universal
Borcherds-weight identity $\kBKM(\Phi_N)=c_N(0)/2$ at $N\in\{1,2,3,4,6\}$;
this note ($N=1$) gives $c_1(0)/2=10/2=5$, matching $\mathrm{wt}(\Delta_5)
=5$.
\item \texttt{chapters/examples/cy\_d\_kappa\_stratification.tex}
Thm.~\texttt{thm:borcherds-weight-kappa-BKM-universal}: the
universal-$N$ statement whose $N=1$ instance is inscribed here.
\end{itemize}

\end{document}
